% Chapter 1: Correlation, Association, and the Yule–Simpson Paradox
% A First Course in Causal Inference - Peng Ding

\chapter{相关性、关联与 Yule-Simpson 悖论}\label{ch:1}

\section{本章导论}

本章我们要回答一个看似简单却至关重要的问题：\textbf{相关性能否直接告诉我们因果关系？}

在日常生活中，我们经常听到这样的话：
\begin{itemize}
  \item "吸烟会导致肺癌"
  \item "教育水平越高，收入越高"
  \item "经常锻炼的人更长寿"
\end{itemize}

这些陈述都试图建立因果关系。但我们是如何得出这些结论的？仅仅因为我们观察到了相关性（correlation）就足够了么？

\section{为什么这个问题很重要}

想象一下以下场景：

\begin{Example}
一家公司想要知道某种新药是否有效。他们给一组患者服用新药，另一组服用安慰剂。结果显示服药组的康复率更高。这能否说明药物有效？
\end{Example}

直觉告诉我们，这似乎能说明问题。但仔细想想，还有很多其他因素可能影响结果：
\begin{itemize}
  \item 也许服药组的患者本身就更健康
  \item 也许医生在给药时，无意中对服药组患者更加关照
  \item 也许这只是随机波动
\end{itemize}

这就是因果推断（causal inference）的核心挑战：\textbf{我们观察到的关联（association），并不等于因果关系（causal relationship）}。

\section{相关性与因果性的区别}

统计学中，我们用不同的术语来区分这些概念：

\begin{itemize}
  \item \textbf{相关性（Correlation）}: 两个变量共同变化的趋势
  \item \textbf{关联（Association）}: 统计学意义上的依赖关系
  \item \textbf{因果性（Causation）}: 一个变量的变化直接导致另一个变量变化
\end{itemize}

关键洞察：\textit{相关性可以是因果的，也可以是非因果的。}

\section{二维列联表与关联度量}

将 $Z$ 视为治疗或暴露，$Y$ 视为结果，可以定义以下关联度量：

\[
\begin{tabular}{lcc}
\hline
        & $Y =1$ & $Y =0$\\\hline
$Z =1$ & $p_{11}$ & $p_{10}$\\
$Z =0$ & $p_{01}$ & $p_{00}$\\
\hline
\end{tabular}
\]

\textbf{风险差（Risk Difference）}：
\[
\RD = \mathbb{P}(Y=1 \mid Z=1) - \mathbb{P}(Y=1 \mid Z=0)
= \frac{p_{11}}{p_{11} + p_{10}} - \frac{p_{01}}{p_{01} + p_{00}}
\]

\textbf{风险比（Risk Ratio）}：
\[
\RR = \frac{\mathbb{P}(Y=1 \mid Z=1)}{\mathbb{P}(Y=1 \mid Z=0)}
= \frac{p_{11}}{p_{11} + p_{10}} \Big / \frac{p_{01}}{p_{01} + p_{00}}
\]

\textbf{优势比（Odds Ratio）}：
\[
\OR = \frac{\mathbb{P}(Y=1 \mid Z=1)/\mathbb{P}(Y=0 \mid Z=1)}{\mathbb{P}(Y=1 \mid Z=0)/\mathbb{P}(Y=0 \mid Z=0)}
= \frac{p_{11} p_{00}}{p_{10} p_{01}}
\]

\textbf{关于记号}：本书使用符号 $\Perp$ 表示随机变量的独立性或条件独立性。

\begin{Proposition}[二维列联表中的独立性]\label{prop:2x2-independence}
\begin{enumerate}
  \item 以下命题等价：$Z \Perp Y$、$\RD = 0$、$\RR = 1$、$\OR = 1$。
  \item 若所有 $p_{zy} > 0$，则 $\RD > 0$ 等价于 $\RR > 1$，也等价于 $\OR > 1$。
  \item 若 $\mathbb{P}(Y=1 \mid Z=1)$ 和 $\mathbb{P}(Y=1 \mid Z=0)$ 都很小，则 $\OR \approx \RR$。
\end{enumerate}
\end{Proposition}

\section{Yule-Simpson 悖论}

悖论的核心思想：在总体中观察到的关联方向，可能与分层后各组中的关联方向相反。

\begin{Example}[经典的 Yule-Simpson 悖论例子]
考虑一个大学录取数据：

\begin{center}
\begin{tabular}{c|c|c}
\hline
院系 & 男生录取率 & 女生录取率 \\
\hline
A & 62\% (341/550) & 82\% (137/167) \\
B & 63\% (28/44) & 68\% (24/35) \\
\hline
\textbf{总体} & \textbf{45\% (369/816)} & \textbf{30\% (161/525)} \\
\hline
\end{tabular}
\end{center}

有趣的现象：
\begin{itemize}
  \item 在 A 院系，男女生录取率分别是 62\% 和 82\%
  \item 在 B 院系，男女生录取率分别是 63\% 和 68\%
  \item 但总体来看，男生的录取率反而更高！
\end{itemize}

这就是 Yule-Simpson 悖论：\textbf{在每个分层中都有利的因素，在总体中可能不利}。
\end{Example}

这个例子揭示了一个深刻的问题：\textit{如果我们不控制混杂因素（confounding factors），就可能得出完全错误的因果结论}。

\section{偏相关系数与 Yule-Simpson 悖论}

考虑三维正态随机向量：
\[
\begin{pmatrix}
X \\ Y \\ Z
\end{pmatrix}
\sim N\left(
\begin{pmatrix}
0 \\ 0 \\ 0
\end{pmatrix},
\begin{pmatrix}
1 & \rho_{XY} & \rho_{XZ} \\
\rho_{XY} & 1 & \rho_{YZ} \\
\rho_{XZ} & \rho_{YZ} & 1
\end{pmatrix}
\right)
\]

$Y$ 和 $Z$ 的偏相关系数 $\rho_{YZ|X}$ 定义为在条件分布 $(Y, Z) \mid X$ 中的相关系数。

\begin{Proposition}[偏相关系数公式]\label{prop:partial-correlation}
\[
\rho_{YZ|X} = \frac{\rho_{YZ} - \rho_{YX}\rho_{ZX}}{\sqrt{1-\rho_{YX}^2}\sqrt{1-\rho_{ZX}^2}}
\]
\end{Proposition}

\textbf{重要含义}：即使 $\rho_{YZ} > 0$，偏相关系数 $\rho_{YZ|X}$ 仍可能为负！这正是正态随机向量情形下的 Yule-Simpson 悖论。

\section{LaLonde 数据与规格搜索问题}

LaLonde 数据集包含参与职业培训项目（National Supported Work Demonstration）的参与者的结果。在观察性研究中，协变量选择的敏感性是一个重要问题。

当我们运行 $2^{10} = 1024$ 种可能的协变量子集回归时，treatment 系数的显著性可能对模型规格高度敏感。这提醒我们：在观察性研究中，规格搜索可能导致误导性结论。

\section{种族歧视的简历实验}

Bertrand and Mullainathan (2004) 通过随机改变简历上的名字来研究种族歧视。这些名字被设计成听起来像非裔美国人或白人的名字。实验显示，有黑人名字的简历回调率显著较低。

当我们分别对男性和女性进行分析时，可能会发现有趣的亚组差异——这进一步说明"平均效应"可能掩盖了重要的异质性。

\section{本章小结与下节预告}

本章的核心要点：
\begin{enumerate}
  \item \textbf{相关性 $\neq$ 因果性}：观察到的关联可能是由于混杂因素造成的
  \item \textbf{关联度量}：$\RD$、$\RR$、$\OR$ 及它们与独立性的关系
  \item \textbf{Yule-Simpson 悖论}：总体关联可能与分层关联方向相反
  \item \textbf{偏相关}：即使边际相关为正，偏相关仍可能为负
  \item \textbf{因果推断的必要性}：我们需要系统性的方法来区分真实的因果效应
\end{enumerate}

下一章我们将学习因果推断的基石——\textbf{潜在结果框架（Potential Outcomes Framework）}，这将为我们的因果分析提供坚实的数学基础。

\section{习题}\label{sec:1-homework}

\begin{HomeworkProblem}[Independence in two-by-two tables]\label{exr:1-1}

Prove (1) and (2) in Proposition \ref{prop::2x2-independence}.


\end{HomeworkProblem}

\begin{HomeworkProblem}[More examples of the Yule--Simpson Paradox]\label{exr:1-2}


Give a numeric example of a two-by-two-by-two table in which the Yule--Simpson Paradox arises.

Find a real-life example in which the Yule--Simpson Paradox arises.



\end{HomeworkProblem}

\begin{HomeworkProblem}[Correlation and partial correlation]\label{exr:1-3}

Consider a three-dimensional Normal random vector:
\[
\begin{pmatrix}
X\\
Y\\
Z
\end{pmatrix}
\sim
\textsc{N}
\left(
\begin{pmatrix}
0\\
0\\
0
\end{pmatrix},
\begin{pmatrix}
1 & \rho_{XY} &\rho_{XZ} \\
  \rho_{XY}  & 1 & \rho_{YZ} \\
  \rho_{XZ} &  \rho_{YZ} & 1
\end{pmatrix}
\right) .
\]
The correlation coefficient between $Y$ and $Z$ is $\rho_{YZ}$.
There are many equivalent definitions of the partial correlation coefficient.
For a multivariate Normal vector, let $\rho_{YZ|X}$ denote the partial correlation coefficient between $Y$ and $Z$ given $X$, which is defined as their correlation coefficient in the conditional distribution $(Y, Z)\mid X$. Show that
\[
\rho_{YZ\mid X} = \frac{   \rho_{YZ} - \rho_{YX} \rho_{ZX}  }{  \sqrt{  1-\rho_{YX}^2 }  \sqrt{1- \rho_{ZX}^2  }  }  .
\]

Give a numerical example with $\rho_{YZ} > 0$ and $\rho_{YZ|X} < 0$.

Remark: This is the Yule--Simpson Paradox for a Normal random vector.
You can use the results in Chapter \ref{sec::multivariate-normal} to prove the formula for the partial correlation coefficient.



\end{HomeworkProblem}

\begin{HomeworkProblem}[Specification searches]\label{exr:1-4}

Section \ref{section:correlatioandregressionintro} re-analyzes the data used by \citet{hainmueller2012entropy}. We used the outcome named \ri{re78} and the treatment named \ri{treat} in the analysis. Moreover, the data also contain $10$ covariates and therefore $2^{10} = 1024$ possible subsets of covariates in the linear regression. Run $1024$ linear regressions with all possible subsets of covariates, and report the regression coefficients of the treatment. How many coefficients of the treatment are positively significant, how many are negatively significant, and how many are not significant? You can also report other interesting findings from these regressions.


\end{HomeworkProblem}

\begin{HomeworkProblem}[More on racial discrimination]\label{exr:1-5}

Section \ref{sec::twobytwotables-withresume} re-analyzes the data collected by \citet{bertrand2004emily}. Conduct analyses separately for males and females. What do you find from these subgroup analyses?


\end{HomeworkProblem}

\begin{HomeworkProblem}[Recommended reading]\label{exr:1-6}

\citet{bickel1975sex} is the original paper for the paradox reported in Section \ref{sec::berkeleydata}.
\citet{pearl2018book} recently revisited the study from the causal inference perspective.



\end{HomeworkProblem}
